%
% Chapter — Conclusion
%
\chapter{Conclusion} \label{cap:conclusion}

This report has presented the design and implementation of \textit{Diário LX}, a
purpose-built digital journalism platform developed in collaboration with the Escola
Superior de Comunicação Social (ESCS) and the LIACOM research unit. The project set out
to support the complete editorial lifecycle --- from content authoring to publication
--- through a role-based workflow, a structured multimedia content model, and a
dedicated public website, within a decoupled client--server architecture. As discussed
in Chapter~\ref{cap:background}, existing content management systems such as WordPress
and Ghost were evaluated but did not fully accommodate these requirements, motivating
the development of a tailored solution. This document describes the final version of
that system, which has been implemented and deployed.

\section{Summary of Work} \label{sec:conc_summary}

The delivered platform realises the architecture and requirements described in the
preceding chapters:

\begin{itemize}
      \item \textbf{Editorial workflow}: the full content lifecycle --- draft, pending
            approval, published, and rejected --- is implemented and enforced according
            to the actor's role. Contributors submit content for review; Editors and
            Administrators approve it (publishing it) or reject it with a mandatory
            comment, and every transition is recorded in the content's history.

      \item \textbf{Content model and editor}: the backoffice content editor supports
            block-based authoring across all content types --- articles, videos,
            podcasts, episodes, and photo essays --- with inline formatting, single media,
            captioned image galleries, external embeds, and full metadata configuration.
            Media assets carry credits and accessible descriptions.

      \item \textbf{Backoffice}: the editorial interface covers content management, user
            administration, invitation-based onboarding, category and tag management, the
            media library, and curation of the public homepage.

      \item \textbf{Public website}: the reader-facing site presents a homepage driven by
            the backoffice curation, article pages, category, tag, and section listings,
            dedicated podcast, episode, and photo-essay pages, and team and author pages,
            with a responsive layout.

      \item \textbf{Security}: authentication uses short-lived JWT access tokens and
            long-lived opaque refresh tokens, issued as \texttt{HttpOnly},
            \texttt{SameSite=Strict}, and --- in production --- \texttt{Secure} cookies.
            Role-based access control is enforced at the HTTP boundary.

      \item \textbf{Deployment}: the system is fully containerised and its images are
            published to a container registry through continuous integration, allowing
            the production server to deploy by pulling rather than building, behind a
            TLS-terminating reverse proxy (Chapter~\ref{cap:deployment}).

      \item \textbf{Quality assurance}: automated backend and end-to-end frontend tests
            run in a continuous integration pipeline, and the REST API is documented
            through an interactive OpenAPI interface.
\end{itemize}

\section{Final Remarks} \label{sec:conc_remarks}

\textit{Diário LX} demonstrates the feasibility of a purpose-built solution tailored to a
specific editorial context, where a general-purpose content management system would have
imposed structural and workflow compromises. The resulting platform delivers a complete,
deployed system that covers the editorial process end to end, together with the public
experience through which its content is read. Beyond the product itself, the project
provided practical experience in full-stack development, layered software architecture,
functional error handling, and containerised deployment, applied to the concrete
requirements of a real editorial partner.
